\documentclass{article}
\usepackage{tcee-repaso}
\tcetema[Mercado único y competencia en la UE]{3.B.39}{El mercado único de la Unión Europea. El principio de libre circulación de mercancías, servicios, personas y capitales. La política de competencia.}
\tcefecha{19/11/2026}
\begin{document}
\cabecerarepaso
\begin{columnas}

% ==========================================================================
\bloque[2,5]{Introducción}
\begin{itemize}
\item \textbf{Enganche}: «El pez grande se come al chico».
  \begin{itemize}
  \item \autor{Cecchini} (1988): en más de la mitad de las ramas cabían 20 empresas eficientes en el mercado comunitario y solo 4 en el mayor nacional.
  \item \autor{Draghi} (2024): las empresas europeas siguen siendo pequeñas; nuevas directrices de concentraciones (2026).
  \end{itemize}
\item \textbf{Relevancia}: 450 M de consumidores; +3-4\,\% de PIB (Comisión); barreras = aranceles del 44\,\% (bienes) y 110\,\% (servicios) (FMI).
  \begin{itemize}
  \item España: más de la mitad de sus exportaciones de bienes \vertema{4.A.22}.
  \end{itemize}
\item \textbf{Contextualización}: \concepto{mercado interior} (art.~26 TFUE): espacio sin fronteras interiores con las 4 libertades.
  \begin{itemize}
  \item Integración \emph{negativa} y \emph{positiva} (\autor{Tinbergen}, 1954).
  \item Mercado interior, competencia compartida; normas de competencia, exclusiva \vertema{3.B.36}.
  \end{itemize}
\item \textbf{Problemática}: ¿por qué conviene y cómo se construyó? ¿Qué falta en cada libertad? ¿Por qué una competencia común? ¿Campeones?
\item \textbf{Estructura}: \hyperref[rep:I]{I} el mercado; \hyperref[rep:II]{II} las libertades; \hyperref[rep:III]{III} la competencia.
\end{itemize}
\esquemacante{%
\apartadocante{Introducción}{2,5}
\apartadocante{I. El mercado único}{7}
\subcante{1. Justificación: competencia y escala}{3}
\subcante{2. Fundamentos jurídicos}{1,5}
\subcante{3. De Roma a la estrategia de 2025}{2,5}
\apartadocante{II. Las 4 libertades}{9}
\subcante{1. Bienes}{3}
\subcante{2. Servicios}{2,5}
\subcante{3. Personas}{1,5}
\subcante{4. Capitales}{2}
\apartadocante{III. La política de competencia}{10}
\subcante{1. Motivación}{1}
\subcante{2. Marco jurídico: los 4 pilares}{8}
\subcante{3. Promoción, coordinación y debate}{1}
\apartadocante{Conclusión}{1,5}}

% ==========================================================================
\bloque[7]{I. El mercado único\label{rep:I}}
\begin{ideaclave}
Se hizo en 2 tiempos (unión aduanera de 1968 y programa de 1992); su ganancia es la competencia y la escala: mercados más grandes, márgenes más pequeños y empresas más eficientes. Sigue incompleto, sobre todo en servicios.
\end{ideaclave}

\sub{I.1. La justificación económica: competencia y escala\label{rep:I.1}}
\begin{itemize}
\item \emph{Teorías}: ventaja comparativa \vertema{3.B.5}; escala, variedad y heterogeneidad \vertema{3.B.6}; integración \vertema{3.B.10}; crecimiento \vertema{3.B.9}.
\end{itemize}
\subsub{a) La integración de mercados oligopolísticos}
\begin{itemize}
\item \emph{Idea} (\autor{Smith}-\autor{Venables}, 1988):
  \begin{itemize}
  \item más rivales, demanda más elástica: \emph{efecto procompetitivo};
  \item salen empresas y las demás crecen: \emph{racionalización}.
  \end{itemize}
\item \emph{Resultado}, con libre entrada:
\funcionclave{p^{*}-c=\sqrt{\frac{\mathrm{CF}}{S}}}
  \begin{itemize}
  \item $S\to2S$: margen $-$29\,\%, producción por empresa +41\,\%, menos empresas.
  \end{itemize}
\item \emph{Implicaciones}: la integración trae concentración $\Rightarrow$ hace falta competencia; las barreras no arancelarias mantienen $S$ pequeño.
\end{itemize}
\subsub{b) La evidencia}
\begin{itemize}
\item Comercio intra-UE de bienes 17\,\% del PIB (2002) $\to$ 22\,\% (2025); servicios 4,8 $\to$ 8\,\%.
\item Márgenes: bajan en industria, no en servicios (\autor{Badinger}, 2007).
\item PIB: +3-4\,\% (Comisión) a $\approx$9\,\% (\autor{in 't Veld}, 2019); 4,4\,\% de bienestar (\autor{Mayer} et al., 2019).
\item FMI: 44/110\,\%; quitarlas, +7\,\% de productividad.
\end{itemize}
\graficorepaso[0.78\linewidth]{comercio-intra-ue}{Comercio intra-UE (\% del PIB), 2002-2025}

\sub{I.2. Fundamentos jurídicos del mercado interior\label{rep:I.2}}
\begin{itemize}
\item \emph{Art.~26 TFUE} (del art.~8A del Acta Única):
  \begin{itemize}
  \item mandato: la Unión adopta las medidas;
  \item definición: «espacio sin fronteras interiores» con las 4 libertades;
  \item progreso equilibrado (Consejo).
  \end{itemize}
\item Compartida (art.~4.2.a); subsidiariedad y proporcionalidad; competencia y unión aduanera, \emph{exclusivas} (3.1.a y b).
\item \emph{Art.~114}: PLO, mayoría cualificada; nivel de protección elevado; 114.4-6, normas nacionales más protectoras.
  \begin{itemize}
  \item Excluye fiscalidad (113 y 115: unanimidad), personas y trabajadores. Libertades: arts.~28-66.
  \end{itemize}
\end{itemize}
\graficorepaso[0.92\linewidth]{base-juridica}{Del art.~26 al 114 y a cada libertad}

\sub{I.3. De Roma a la estrategia de 2025\label{rep:I.3}}
\subsub{a) La unión aduanera y el estancamiento (1957-1985)}
\begin{itemize}
\item Roma (1957): mercado común. 1-7-1968: unión aduanera (arancel exterior común).
\item 1973-1985: barreras no arancelarias y unanimidad; el TJ resiste: \emph{Dassonville} (1974), \emph{Cassis} (1979).
\end{itemize}
\subsub{b) El programa de 1992}
\begin{itemize}
\item Libro Blanco (1985, \autor{Delors}, \autor{Cockfield}): $\approx$280 medidas contra barreras físicas, técnicas y fiscales.
\item \emph{Acta Única} (1986; vigor 1987): mercado interior para 1992 y mayoría cualificada.
\item \autor{Cecchini} (1988): +4,5\,\% PIB, $-$6,1\,\% precios, +1,8 M de empleos ($\pm$30\,\%).
  \begin{itemize}
  \item Fronteras: 1,8\,\% del comercio; escala: 2\,\% del PIB.
  \end{itemize}
\item 1-1-1993: mercado único.
\end{itemize}
\subsub{c) Los relanzamientos (2010-2026)}
\begin{itemize}
\item \autor{Monti} (2010); mercado único digital y UMC (2015).
\item \autor{Letta} (2024): 5.ª libertad, Unión de Ahorros e Inversiones, 28.º régimen.
\item \autor{Draghi} (2024): 750.000-800.000 M€/año (4,4-4,7\,\% del PIB).
\item IMERA (Reg. 2024/2747; aplicable 29-5-2026): emergencias del mercado interior.
\item \emph{Estrategia} (21-5-2025): las \emph{Terrible Ten} barreras; hoja de ruta «Una Europa, un mercado» (2026-2027).
\item EU Inc. (18-3-2026): 28.º régimen, 48 h, $<$100\,€, sin capital mínimo.
\end{itemize}
\graficorepaso[0.95\linewidth]{linea-tiempo}{De Roma (1957) a la estrategia de 2025}

% ==========================================================================
\bloque[9]{II. Las 4 libertades de circulación\label{rep:II}}
\begin{ideaclave}
Bienes, la más avanzada; capitales, la más amplia (\emph{erga omnes}); servicios, la más atrasada; personas, la más limitada en la práctica.
\end{ideaclave}
\graficorepaso[0.95\linewidth]{sintesis-libertades}{Mismo patrón, distinto avance}

\sub{II.1. Mercado interior de bienes\label{rep:II.1}}
\subsub{a) Marco jurídico y jurisprudencia}
\begin{itemize}
\item Arts.~28-37: 30 (exacciones), 34-35 (restricciones y medidas de efecto equivalente), 110 (discriminación fiscal).
  \begin{itemize}
  \item Excepciones: art.~36 y exigencias imperativas.
  \end{itemize}
\item \emph{Dassonville} (1974): definición amplia.
\item \emph{Cassis de Dijon} (1979): \concepto{reconocimiento mutuo}.
\item \emph{Keck} (1993): modalidades de venta, fuera.
\item \emph{Fresas españolas} (1997): el Estado debe actuar contra particulares.
\end{itemize}
\subsub{b) Las barreras de Cecchini y su desmontaje}
\subsubsub{Físicas}
\begin{itemize}
\item Fin de controles (1993); Intrastat; IVA sin frontera.
\item Código Aduanero 2026/2108; Autoridad en Lille; envíos $<$150\,€: 3\,€ por artículo (1-7-2026) + 2\,€ de gestión.
\end{itemize}
\subsubsub{Técnicas}
\begin{itemize}
\item Enfoque tradicional (detallado, unanimidad) $\to$ \emph{nuevo enfoque} (1985):
  \begin{itemize}
  \item requisitos esenciales; normas armonizadas voluntarias; presunción de conformidad; marcado CE.
  \end{itemize}
\item Riesgos del reconocimiento mutuo: deslocalización y carrera a la baja.
\item NLF (2008); Reg. 2019/515; TRIS (3 meses); geobloqueo (2018/302); pasaporte digital (2026).
\end{itemize}
\subsubsub{Fiscales}
\begin{itemize}
\item Unanimidad (113); origen frente a destino; estrategia de \emph{aproximación}.
\item IVA: Sexta Directiva (1977); \emph{régimen transitorio} (1993, destino, VIES).
  \begin{itemize}
  \item Dir. 2022/542: general $\geq$15\,\%, 2 reducidos $\geq$5\,\%, uno $<$5\,\% y tipo cero; España 21/10/4 \vertema{4.B.20}.
  \item ViDA (2025/516); brecha del IVA 128.000 M€ (9,5\,\%, 2023); fraude carrusel.
  \end{itemize}
\end{itemize}
\subsubsub{Contratos públicos}
\begin{itemize}
\item 15\,\% del PIB (2,6 billones); directivas de 2014.
\end{itemize}
\subsub{c) Valoración}
\begin{itemize}
\item (+) Comercio, convergencia de precios, integración del Este.
\item ($-$) Estancamiento en la UE-15, \emph{home bias}; el 60\,\% de las barreras, igual que hace 20 años.
\end{itemize}

\sub{II.2. Mercado interior de servicios\label{rep:II.2}}
\begin{itemize}
\item 75\,\% del PIB, 8\,\% del PIB comerciado; en valor añadido, $\approx$1/2 del comercio \vertema{3.B.34}.
\item \emph{Marco}: establecimiento (art.~49; modo 3) y prestación (56; modos 1, 2 y 4).
  \begin{itemize}
  \item Excepciones: 51 (poder público), 52 y 62 (orden, seguridad, salud).
  \item \emph{Van Binsbergen} (1974); \emph{Laval} y \emph{Viking} (2007).
  \end{itemize}
\item \emph{Implementación}:
  \begin{itemize}
  \item años 90: sectorial (redes);
  \item \concepto{Directiva de servicios} (2006/123; Leyes 17 y 25/2009): autorizaciones necesarias y proporcionadas, ventanillas únicas, IMI; sin país de origen;
  \item cualificaciones (2005/36): automático en 7 profesiones; test de proporcionalidad (2018/958);
  \item pasaporte financiero \vertema{3.B.44}; DSA y DMA.
  \end{itemize}
\item \emph{Retos}: arancel equivalente $\leq$110\,\%; telecomunicaciones: oligopolio sin escala.
\end{itemize}

\sub{II.3. Libre circulación de personas\label{rep:II.3}}
\begin{itemize}
\item \emph{Justificación}: neoclásica, $w/r$ igual en todas partes; geografía: aglomeración. \autor{Card} (1990): Mariel, efecto pequeño.
\item \emph{Marco}: art.~45 (límites: orden, seguridad, salud; empleo público); Reg. 492/2011.
  \begin{itemize}
  \item Ciudadanía (art.~21); Dir. 2004/38: 3 meses; más (recursos y seguro); permanente a los 5 años.
  \end{itemize}
\item \emph{Schengen} (1985/1990): 29 miembros (Bulgaria y Rumanía desde 2025; fuera Irlanda y Chipre).
  \begin{itemize}
  \item Controles temporales hasta 2 + 1 años (Código 2024); EES pleno (abril de 2026).
  \end{itemize}
\item \emph{Desplazados}: 96/71 $\to$ 2018/957 (misma remuneración; a los 12-18 meses, Derecho laboral de acogida).
\item 14 M viven o trabajan en otro Estado; movilidad baja \vertema{3.B.16}.
\end{itemize}
\graficorepaso[0.7\linewidth]{mapa-schengen}{El espacio Schengen (nov. 2026)}

\sub{II.4. Libre circulación de capitales\label{rep:II.4}}
\begin{itemize}
\item \emph{Justificación}: intertemporal \vertema{3.B.12}, cartera e IED \vertema{3.B.18}; imprescindible en la UEM.
\item \emph{Marco}: \concepto{Dir. 88/361} (1-7-1990; España, Irlanda, Grecia, Portugal hasta 1992); Maastricht (1994).
  \begin{itemize}
  \item Art.~63 (\emph{erga omnes}), 64, 65 (fiscal, prudencial, orden público: Chipre 2013, Grecia 2015), 66 (6 meses).
  \item Acciones de oro: C-463/00 (2003).
  \end{itemize}
\item \emph{IED}: 2019/452 $\to$ 2026/1386 (control obligatorio) \vertema{3.B.41}. UMC $\to$ UAI (2025) \vertema{3.B.44}.
\item \emph{Fiscalidad directa} (115, unanimidad) \vertema{4.B.18}:
  \begin{itemize}
  \item matrices y filiales; código de conducta (1997); ATAD (2016);
  \item pilar 2 (2022/2523, 15\,\%); BEFIT sin acuerdo; IS del 9\,\% (HU) al 30\,\% (DE);
  \item intercambio de información: DAC (2011/16; CRS 2016).
  \end{itemize}
\item \emph{Evidencia}: \autor{Feldstein}-\autor{Horioka} (1980): ahorro e inversión correlacionados, menos en la zona del euro; \autor{Gopinath} et al. (2017): capital al Sur mal asignado.
\end{itemize}

% ==========================================================================
\bloque[10]{III. La política europea de competencia\label{rep:III}}
\begin{ideaclave}
Es la otra cara del mercado único: controla a las empresas (acuerdos, abusos, concentraciones) y a los Estados (ayudas); criterio, el bienestar del consumidor; debate: escala y resiliencia.
\end{ideaclave}

\sub{III.1. Motivación de la teoría económica\label{rep:III.1}}
\begin{itemize}
\item Integración positiva: la ganancia de \hyperref[rep:I.1]{I.1} exige competencia.
\item Eficiencia asignativa, productiva (ineficiencia X, \autor{Leibenstein}) y dinámica.
\item Harvard, Chicago (\autor{Stigler}), Friburgo; enfoque económico.
\end{itemize}

\sub{III.2. Marco jurídico: los 4 pilares\label{rep:III.2}}
\begin{itemize}
\item Exclusiva (art.~3.1.b); arts.~101-109, Reg. 1/2003 y 139/2004; + subvenciones extranjeras y DMA.
\item \emph{Quién}: Colegio (DG COMP, \autor{Ribera}), ANC (CNMC) y jueces; TG y TJ.
  \begin{itemize}
  \item Autoridad «mejor situada»; la Comisión si $>$3 Estados.
  \end{itemize}
\item Evolución: Reg. 17/1962 (notificación) $\to$ 1/2003 y 139/2004 $\to$ DMA y subvenciones extranjeras (2022).
\end{itemize}
\graficorepaso[0.92\linewidth]{pilares-competencia}{4 pilares clásicos y 2 instrumentos nuevos; quién aplica y controla}
\subsub{a) Primer pilar: los acuerdos entre empresas (art.~101)}
\begin{itemize}
\item Acuerdos, decisiones, prácticas concertadas: nulos.
\item \emph{101.3}: 4 condiciones (mejora, participación de los usuarios, indispensables, no eliminan la competencia).
  \begin{itemize}
  \item Exenciones por categorías: verticales 2022/720, horizontales 2023 (sostenibilidad); \emph{de minimis} 10/15\,\%.
  \end{itemize}
\item \emph{Reg. 1/2003}: excepción legal; Red Europea; multas $\leq$10\,\%.
\item Clemencia: exención, 30-50, 20-30, $\leq$20\,\%; transacción, 10\,\%.
\end{itemize}
\subsubsub{El coste del cártel y la multa disuasoria}
\begin{itemize}
\item \autor{Becker} (1968), \autor{Landes} (1983): sobreprecio $\Delta B=(p^K-c)Q^K$ y PIE (\autor{Harberger}).
\item Se forma si $\Delta B>\pi M$:
\funcionclave{M\geq\frac{\Delta B}{\pi}}
  \begin{itemize}
  \item $\pi=1/6\Rightarrow M=6\Delta B$.
  \item Clemencia \vertema{3.A.19}: baja el beneficio de coludir y sube el de desviarse: hace falta más paciencia.
  \end{itemize}
\item \emph{Evidencia}: sobreprecio mediano 23\,\% (\autor{Connor}); $\pi\approx13$\,\%.
  \begin{itemize}
  \item Multas 30.600 M€, 162 decisiones (1990-2025); camiones 3.807. El tope del 10\,\% disuade poco.
  \end{itemize}
\end{itemize}
\pareja{cartel-bienestar}{Sobreprecio $\Delta B$ y PIE}{multas-carteles}{Multas por cártel (M€), 1990-2025}
\subsub{b) Segundo pilar: el abuso de posición dominante (art.~102)}
\begin{itemize}
\item Se prohíbe el abuso, no el dominio.
  \begin{itemize}
  \item Mercado relevante (comunicación de 2024) \vertema{3.A.20}; dominio $>$50\,\% (presunción), $<$40\,\% improbable.
  \item Abusos excluyentes y explotadores; directrices de 3-9-2026.
  \end{itemize}
\item Google Shopping 2.424 M€; Android 4.125 (TJ, 2026); Apple 1.840; Meta 798; adtech 2.950.
\item \emph{DMA} (\emph{ex ante}): 7 guardianes, 23 servicios; Apple 500, Meta 200 (2025), Google 890 (2026).
\end{itemize}
\subsub{c) Tercer pilar: el control de concentraciones}
\begin{itemize}
\item Reg. 139/2004 (test SIEC): 5.000/250 M€ o 2.500/100/25/100; regla de los 2/3.
\item Fase I, 25 días; fase II, 90. 10.200 notificaciones, 33 prohibiciones.
  \begin{itemize}
  \item \emph{Siemens}/\emph{Alstom} (2019), \emph{Booking}/\emph{eTraveli} (2023); \emph{Illumina} (2024): sin remisiones por debajo de umbrales.
  \end{itemize}
\end{itemize}
\subsubsub{El compromiso de \autor{Williamson} (1968)}
\begin{itemize}
\item PIE $A_1\approx\tfrac12\Delta p\,\Delta Q$ frente a ahorro de costes $A_2=\Delta c\,Q_1$:
\funcionclave{\frac{\Delta c}{c_0}>\frac{|\varepsilon|}{2}\left(\frac{\Delta p}{p_0}\right)^{2}}
  \begin{itemize}
  \item $|\varepsilon|=1$, +10\,\% de precio: basta $-$0,5\,\% de coste.
  \end{itemize}
\item UE: bienestar del consumidor (eficiencias trasladadas).
\item \autor{Farrell}-\autor{Shapiro} (1990); \autor{Kwoka}: +7\,\% de precio; adquisiciones asesinas.
\item Borrador de directrices (30-4-2026): escala, resiliencia, innovación.
\end{itemize}
\pareja{williamson}{$A_1$ pequeño frente a $A_2$}{concentraciones}{Notificaciones y fases II, 1991-2025}
\subsub{d) Cuarto pilar: las ayudas de Estado}
\begin{itemize}
\item \emph{Art.~107.1}: recursos estatales, ventaja, selectividad, efecto sobre el comercio.
  \begin{itemize}
  \item 107.2, automáticas; 107.3: regiones, IPCEI (11), actividades, cultura.
  \item 108: notificación y suspensión; recuperación.
  \end{itemize}
\item RGEC (651/2014); \emph{de minimis} 300.000\,€/3 años; marcos COVID, Ucrania, CISAF (2025-2030).
\item 2024: 168.230 M€, 0,94\,\% del PIB (2020: 2,41\,\%); España 0,99\,\%; Alemania $\approx$25\,\%.
\item \emph{Carrera de subvenciones}: dilema del prisionero (10,10) $\to$ (7,7); la prohibición es la cooperación (\autor{Collie}, 2000).
\item \emph{Subvenciones extranjeras} (2022/2560): 500/50 M€; contratos 250 M€.
  \begin{itemize}
  \item CRRC, \emph{e\&}/PPF, ADNOC/Covestro; 273 notificaciones.
  \end{itemize}
\end{itemize}
\graficorepaso[0.92\linewidth]{ayudas-estado}{Ayudas de Estado: pico de 2020 y reparto por Estados en 2024}

\sub{III.3. Promoción, coordinación y debate actual\label{rep:III.3}}
\begin{itemize}
\item Industrias de red: telecomunicaciones, energía (4 paquetes), transporte; separar red y servicios \vertema{3.A.20}.
\item ECN+ (2019/1); Red Internacional de Competencia.
\item \emph{¿Campeones europeos?}
  \begin{itemize}
  \item A favor: costes fijos y mercado mundial (\autor{Draghi}).
  \item En contra: precios en casa; \autor{Gutiérrez}-\autor{Philippon} (2023): la fragmentación, no las fusiones, frena el crecimiento.
  \end{itemize}
\end{itemize}

% ==========================================================================
\bloque[1,5]{Conclusión}
\begin{itemize}
\item \textbf{Recapitulación}: 2 tiempos (1968, 1992); competencia y escala; libertades desiguales; competencia frente a empresas y Estados.
\item \textbf{Relevancia}: principal mercado y marco de la competencia de España \vertema{4.A.15}.
\item \textbf{Extensiones}: \vertema{3.B.10} \vertema{3.B.36} \vertema{3.B.40} \vertema{3.B.41} \vertema{3.B.44}.
\item \textbf{Opinión}: completar servicios, capitales y digital antes que fabricar campeones.
  \begin{itemize}
  \item Directrices con innovación y resiliencia, sin abandonar al consumidor.
  \end{itemize}
\item \textbf{Idea final}: el mercado agranda el estanque y la competencia vigila que se crezca por ser mejor.
  \begin{itemize}
  \item La solución no es dejar que el grande se coma al chico, sino quitar las redes que parten el estanque.
  \end{itemize}
\end{itemize}

% ==========================================================================
\anexorepaso
\bloque{Datos clave}
\begin{itemize}
\item Intra-UE: \dato{bienes 21,8\,\%, servicios 8\,\% del PIB}{Eurostat, 2025}; FMI: 44/110\,\%.
\item Cárteles: \dato{30.600 M€}{1990-2025}. Concentraciones: 10.200, 33 prohibiciones.
\item Ayudas: \dato{0,94\,\% del PIB}{2024}; España 0,99\,\%.
\item Schengen: 29; DMA: 7 guardianes.
\end{itemize}
\bloque{Autores y fechas}
\hito{1968}{Unión aduanera; \autor{Williamson}; \autor{Becker}}
\hito{1979}{\emph{Cassis de Dijon}}
\hito{1985}{Libro Blanco; nuevo enfoque}
\hito{1988}{\autor{Cecchini}; \autor{Smith}-\autor{Venables}}
\hito{1993}{Mercado único; \emph{Keck}}
\hito{2006}{Directiva de servicios}
\hito{2022}{DMA, DSA, subvenciones extranjeras}
\hito{2024}{\autor{Letta}, \autor{Draghi}}
\bloque{Preguntas del tribunal}
\begin{enumerate}
\item \textbf{¿Qué aportó \emph{Cassis}?} Reconocimiento mutuo y exigencias imperativas.
\item \textbf{¿Por qué van atrasados los servicios?} Regulación nacional, movilidad del proveedor, Directiva rebajada.
\item \textbf{¿Cuánto debe ser la multa a un cártel?} $\Delta B/\pi$; el tope del 10\,\% limita la disuasión.
\item \textbf{¿Fusiones con eficiencias?} \autor{Williamson}: sí con bienestar total; con el del consumidor, si no sube el precio.
\item \textbf{¿Art.~102 o DMA?} \emph{Ex post}, caso a caso, frente a \emph{ex ante}, obligaciones fijas.
\item \textbf{¿Por qué controlar ayudas nacionales?} Dilema del prisionero entre Estados; asimetría fiscal.
\end{enumerate}
\begin{trampa}
\begin{itemize}
\item Unión aduanera \emph{1-7-1968}; mercado único \emph{1-1-1993}.
\item \emph{Keck} no anula \emph{Cassis}: saca las \emph{modalidades de venta}.
\item 101.3: \emph{4} condiciones acumulativas.
\item Fase I: \emph{25} días hábiles; fase II: \emph{90}.
\item Excepciones de servicios: arts.~\emph{51 y 52} (no el 36).
\item Schengen: \emph{29} miembros; fuera Irlanda y Chipre.
\end{itemize}
\end{trampa}

\end{columnas}
\end{document}
